% !TEX root = ../main.tex

\section{The Quantum Generator}
\label{sec:quantum}
\frame[plain]{\sectionpage}

%----------------------------------------
\begin{frame}{From noise to market scenario: hybrid pipeline}
    \begin{center}
    \begin{tikzpicture}[node distance=6mm,
        box/.style={draw=zhawaccent,thick,rounded corners,
                    minimum width=20mm,minimum height=9mm,
                    fill=zhawaccent!10,align=center,font=\small},
        data/.style={font=\small\itshape},
        arrow/.style={-Latex,thick}]
    \node[data] (z) {$\mathbf{z}\sim\mathcal{N}(0,I_{n_q})$};
    \node[box,right=of z] (Q) {Quantum\\circuit $Q_\theta$};
    \node[data,right=of Q] (mu) {$\mu_i = \langle Z_i\rangle$};
    \node[box,right=of mu] (h) {Linear+\\$\tanh$};
    \node[data,right=of h] (out)
        {$W\!\times\!n_a$ scenario};
    \draw[arrow] (z) -- (Q);
    \draw[arrow] (Q) -- (mu);
    \draw[arrow] (mu) -- (h);
    \draw[arrow] (h) -- (out);
    \end{tikzpicture}
    \end{center}
    \medskip

    \textbf{Two-stage architecture}
    \begin{itemize}
        \item \textbf{Quantum part} $Q_\theta:
        \mathbb{R}^{n_q}\to[-1,1]^{n_q}$ produces $n_q$
        expectation values via an entangled circuit
        \item \textbf{Classical head} upsamples to full output
        dimension via a single \texttt{Linear}~+~$\tanh$
    \end{itemize}
    \medskip

    \emph{Discriminator unchanged from the classical baseline.
    Only the generator differs --- a controlled comparison.}
\end{frame}

%----------------------------------------
\begin{frame}{The variational quantum circuit}
    \begin{columns}[T]
    \begin{column}{.55\textwidth}
        Operations on $\mathcal{H} = \mathbb{C}^{2^{n_q}}$:
        \medskip

        \begin{enumerate}
            \item \textbf{Initial state} $|\mathbf{0}\rangle$
            \item \textbf{AngleEmbedding}: apply $RY(z_i)$ on
            each qubit (encodes input as rotations)
            \item \textbf{Variational layer (×$L$)}:
            trainable $RY(\theta_{\ell,i})$ on each qubit, then
            a ring of CNOTs
            \item \textbf{Measure} $\mu_i = \langle Z_i\rangle
            \in [-1,1]$ on each qubit
        \end{enumerate}
        \medskip

        Total trainable angles: $n_q \times L$.\\
        For $n_q=8, L=3$: \textbf{24 quantum parameters} on a
        256-dim complex state space.
    \end{column}
    \hfill
    \begin{column}{.43\textwidth}
        \begin{block}{Why this design?}
            \begin{itemize}
                \item Hardware-efficient: nearest-neighbour
                CNOTs only
                \item $RY$-only rotations mitigate barren
                plateaus relative to richer
                ansätze~\cite{McClean2018BarrenPlateaus}
                \item Local Pauli-Z measurements bound
                gradient variance~\cite{Cerezo2021LocalCost}
                \item $L=3$: enough depth for entanglement to
                spread; deeper risks barren plateaus
            \end{itemize}
        \end{block}
    \end{column}
    \end{columns}
\end{frame}
